\documentclass{article}
\usepackage{mathtools} 
\usepackage{fontspec}
\usepackage[UTF8]{ctex}
\usepackage{amsthm}
\usepackage{mdframed}
\usepackage{xcolor}
\usepackage{amssymb}
\usepackage{amsmath}


% 定义新的带灰色背景的说明环境 zremark
\newmdtheoremenv[
  backgroundcolor=gray!10,
  % 边框与背景一致，边框线会消失
  linecolor=gray!10
]{zremark}{说明}

% 通用矩阵命令: \flexmatrix{矩阵名}{元素符号}{行数}{列数}
\newcommand{\flexmatrix}[4]{
  \[
  #1 = \begin{pmatrix}
    #2_{11}     & #2_{12}     & \cdots & #2_{1#4}   \\
    #2_{21}     & #2_{22}     & \cdots & #2_{2#4}   \\
    \vdots      & \vdots      & \ddots & \vdots     \\
    #2_{#31}    & #2_{#32}    & \cdots & #2_{#3#4}
  \end{pmatrix}
  \]
}

% 简化版命令（默认矩阵名为A，元素符号为a）: \quickmatrix{行数}{列数}
\newcommand{\quickmatrix}[2]{\flexmatrix{A}{a}{#1}{#2}}

\begin{document}
\title{注释 17.4}
\author{张志聪}
\maketitle

\begin{zremark}
  $f_{xy}(x, y)$在点$(x_0, y_0)$连续，则存在$(x_0, y_0)$的领域$U$，
  对任意$(x, y) \in U$，$f_x$存在。
\end{zremark}
\textbf{证明：}
由连续的定义
\begin{align*}
  \lim\limits_{(x, y) \to (x_0, y_0)} f_{xy}(x, y) = f_{xy}(x_0, y_0)
\end{align*}
可知，存在$(x_0, y_0)$的领域$U$，使得$f_{xy}$在$U$内每一个点都有定义；

对任意$(x_1, y_1) \in U$，由偏导数的定义，我们有
\begin{align*}
  \lim\limits_{\Delta y \to 0} \frac{f_x(x_1, y_1 + \Delta y) - f_x(x_1, y_1)}{\Delta y}
\end{align*}
以上隐含了$f_x(x_1, y_1)$存在，否则上式没有意义。

综上，命题成立。

\begin{zremark}
  命题

  设$z = f(x, y), x = \varphi (s, t), y = \psi (s, t)$，
  若函数$f, \varphi, \psi$都具有连续的二阶偏导数，则$z$对
  $s, t$同样存在二阶连续偏导数。
\end{zremark}

\begin{zremark}
  为什么

  $D$是矩形区域$[a, b] \times [c, d]$，那就不能保证对$D$上的任意两点$P,Q$都有$(8)$是成立。
\end{zremark}
\textbf{证明：}
因为如果$P,Q$都在同一条边上，则$PQ$线段上的点都是边界点，
由书中的“注”可知$(8)$式不成立。

\begin{zremark}
  证明：定理17.8（中值定理）推论。
\end{zremark}

\begin{zremark}
  证明：定理17.10（极值必要条件）
\end{zremark}
\textbf{证明：}
以$f_x(x_0, y_0) = 0$为例。

反证法，假设$f_x(x_0, y_0) \neq 0$，不妨设$f_x > 0$，
由偏导数的定义可知
\begin{align*}
  \lim\limits_{\Delta x \to 0} \frac{f(x_0 + \Delta x, y_0) - f(x_0, y_0)}{\Delta x} 
  = f_x(x_0, y_0) > 0
\end{align*}
由保号性可知，存在$\delta > 0$，$0 < \Delta x < \delta$时，有
\begin{align*}
  \frac{f(x_0 + \Delta x, y_0) - f(x_0, y_0)}{\Delta x} & > 0 \\
  f(x_0 + \Delta x, y_0) - f(x_0, y_0) & > 0 \\
  f(x_0 + \Delta x, y_0) & > f(x_0, y_0)
\end{align*}
这与$f(x_0, y_0)$是极值点矛盾。

\end{document}